\documentclass[11pt]{article}

    \usepackage[breakable]{tcolorbox}
    \usepackage{parskip} % Stop auto-indenting (to mimic markdown behaviour)
    

    % Basic figure setup, for now with no caption control since it's done
    % automatically by Pandoc (which extracts ![](path) syntax from Markdown).
    \usepackage{graphicx}
    % Keep aspect ratio if custom image width or height is specified
    \setkeys{Gin}{keepaspectratio}
    % Maintain compatibility with old templates. Remove in nbconvert 6.0
    \let\Oldincludegraphics\includegraphics
    % Ensure that by default, figures have no caption (until we provide a
    % proper Figure object with a Caption API and a way to capture that
    % in the conversion process - todo).
    \usepackage{caption}
    \DeclareCaptionFormat{nocaption}{}
    \captionsetup{format=nocaption,aboveskip=0pt,belowskip=0pt}

    \usepackage{float}
    \floatplacement{figure}{H} % forces figures to be placed at the correct location
    \usepackage{xcolor} % Allow colors to be defined
    \usepackage{enumerate} % Needed for markdown enumerations to work
    \usepackage{geometry} % Used to adjust the document margins
    \usepackage{amsmath} % Equations
    \usepackage{amssymb} % Equations
    \usepackage{textcomp} % defines textquotesingle
    % Hack from http://tex.stackexchange.com/a/47451/13684:
    \AtBeginDocument{%
        \def\PYZsq{\textquotesingle}% Upright quotes in Pygmentized code
    }
    \usepackage{upquote} % Upright quotes for verbatim code
    \usepackage{eurosym} % defines \euro

    \usepackage{iftex}
    \ifPDFTeX
        \usepackage[T1]{fontenc}
        \IfFileExists{alphabeta.sty}{
              \usepackage{alphabeta}
          }{
              \usepackage[mathletters]{ucs}
              \usepackage[utf8x]{inputenc}
          }
    \else
        \usepackage{fontspec}
        \usepackage{unicode-math}
    \fi

    \usepackage{fancyvrb} % verbatim replacement that allows latex
    \usepackage{grffile} % extends the file name processing of package graphics
                         % to support a larger range
    \makeatletter % fix for old versions of grffile with XeLaTeX
    \@ifpackagelater{grffile}{2019/11/01}
    {
      % Do nothing on new versions
    }
    {
      \def\Gread@@xetex#1{%
        \IfFileExists{"\Gin@base".bb}%
        {\Gread@eps{\Gin@base.bb}}%
        {\Gread@@xetex@aux#1}%
      }
    }
    \makeatother
    \usepackage[Export]{adjustbox} % Used to constrain images to a maximum size
    \adjustboxset{max size={0.9\linewidth}{0.9\paperheight}}

    % The hyperref package gives us a pdf with properly built
    % internal navigation ('pdf bookmarks' for the table of contents,
    % internal cross-reference links, web links for URLs, etc.)
    \usepackage{hyperref}
    % The default LaTeX title has an obnoxious amount of whitespace. By default,
    % titling removes some of it. It also provides customization options.
    \usepackage{titling}
    \usepackage{longtable} % longtable support required by pandoc >1.10
    \usepackage{booktabs}  % table support for pandoc > 1.12.2
    \usepackage{array}     % table support for pandoc >= 2.11.3
    \usepackage{calc}      % table minipage width calculation for pandoc >= 2.11.1
    \usepackage[inline]{enumitem} % IRkernel/repr support (it uses the enumerate* environment)
    \usepackage[normalem]{ulem} % ulem is needed to support strikethroughs (\sout)
                                % normalem makes italics be italics, not underlines
    \usepackage{soul}      % strikethrough (\st) support for pandoc >= 3.0.0
    \usepackage{mathrsfs}
    
	\input{../../macros.tex}
	
    
    % Colors for the hyperref package
    \definecolor{urlcolor}{rgb}{0,.145,.698}
    \definecolor{linkcolor}{rgb}{.71,0.21,0.01}
    \definecolor{citecolor}{rgb}{.12,.54,.11}

    % ANSI colors
    \definecolor{ansi-black}{HTML}{3E424D}
    \definecolor{ansi-black-intense}{HTML}{282C36}
    \definecolor{ansi-red}{HTML}{E75C58}
    \definecolor{ansi-red-intense}{HTML}{B22B31}
    \definecolor{ansi-green}{HTML}{00A250}
    \definecolor{ansi-green-intense}{HTML}{007427}
    \definecolor{ansi-yellow}{HTML}{DDB62B}
    \definecolor{ansi-yellow-intense}{HTML}{B27D12}
    \definecolor{ansi-blue}{HTML}{208FFB}
    \definecolor{ansi-blue-intense}{HTML}{0065CA}
    \definecolor{ansi-magenta}{HTML}{D160C4}
    \definecolor{ansi-magenta-intense}{HTML}{A03196}
    \definecolor{ansi-cyan}{HTML}{60C6C8}
    \definecolor{ansi-cyan-intense}{HTML}{258F8F}
    \definecolor{ansi-white}{HTML}{C5C1B4}
    \definecolor{ansi-white-intense}{HTML}{A1A6B2}
    \definecolor{ansi-default-inverse-fg}{HTML}{FFFFFF}
    \definecolor{ansi-default-inverse-bg}{HTML}{000000}

    % common color for the border for error outputs.
    \definecolor{outerrorbackground}{HTML}{FFDFDF}

    % commands and environments needed by pandoc snippets
    % extracted from the output of `pandoc -s`
    \providecommand{\tightlist}{%
      \setlength{\itemsep}{0pt}\setlength{\parskip}{0pt}}
    \DefineVerbatimEnvironment{Highlighting}{Verbatim}{commandchars=\\\{\}}
    % Add ',fontsize=\small' for more characters per line
    \newenvironment{Shaded}{}{}
    \newcommand{\KeywordTok}[1]{\textcolor[rgb]{0.00,0.44,0.13}{\textbf{{#1}}}}
    \newcommand{\DataTypeTok}[1]{\textcolor[rgb]{0.56,0.13,0.00}{{#1}}}
    \newcommand{\DecValTok}[1]{\textcolor[rgb]{0.25,0.63,0.44}{{#1}}}
    \newcommand{\BaseNTok}[1]{\textcolor[rgb]{0.25,0.63,0.44}{{#1}}}
    \newcommand{\FloatTok}[1]{\textcolor[rgb]{0.25,0.63,0.44}{{#1}}}
    \newcommand{\CharTok}[1]{\textcolor[rgb]{0.25,0.44,0.63}{{#1}}}
    \newcommand{\StringTok}[1]{\textcolor[rgb]{0.25,0.44,0.63}{{#1}}}
    \newcommand{\CommentTok}[1]{\textcolor[rgb]{0.38,0.63,0.69}{\textit{{#1}}}}
    \newcommand{\OtherTok}[1]{\textcolor[rgb]{0.00,0.44,0.13}{{#1}}}
    \newcommand{\AlertTok}[1]{\textcolor[rgb]{1.00,0.00,0.00}{\textbf{{#1}}}}
    \newcommand{\FunctionTok}[1]{\textcolor[rgb]{0.02,0.16,0.49}{{#1}}}
    \newcommand{\RegionMarkerTok}[1]{{#1}}
    \newcommand{\ErrorTok}[1]{\textcolor[rgb]{1.00,0.00,0.00}{\textbf{{#1}}}}
    \newcommand{\NormalTok}[1]{{#1}}

    % Additional commands for more recent versions of Pandoc
    \newcommand{\ConstantTok}[1]{\textcolor[rgb]{0.53,0.00,0.00}{{#1}}}
    \newcommand{\SpecialCharTok}[1]{\textcolor[rgb]{0.25,0.44,0.63}{{#1}}}
    \newcommand{\VerbatimStringTok}[1]{\textcolor[rgb]{0.25,0.44,0.63}{{#1}}}
    \newcommand{\SpecialStringTok}[1]{\textcolor[rgb]{0.73,0.40,0.53}{{#1}}}
    \newcommand{\ImportTok}[1]{{#1}}
    \newcommand{\DocumentationTok}[1]{\textcolor[rgb]{0.73,0.13,0.13}{\textit{{#1}}}}
    \newcommand{\AnnotationTok}[1]{\textcolor[rgb]{0.38,0.63,0.69}{\textbf{\textit{{#1}}}}}
    \newcommand{\CommentVarTok}[1]{\textcolor[rgb]{0.38,0.63,0.69}{\textbf{\textit{{#1}}}}}
    \newcommand{\VariableTok}[1]{\textcolor[rgb]{0.10,0.09,0.49}{{#1}}}
    \newcommand{\ControlFlowTok}[1]{\textcolor[rgb]{0.00,0.44,0.13}{\textbf{{#1}}}}
    \newcommand{\OperatorTok}[1]{\textcolor[rgb]{0.40,0.40,0.40}{{#1}}}
    \newcommand{\BuiltInTok}[1]{{#1}}
    \newcommand{\ExtensionTok}[1]{{#1}}
    \newcommand{\PreprocessorTok}[1]{\textcolor[rgb]{0.74,0.48,0.00}{{#1}}}
    \newcommand{\AttributeTok}[1]{\textcolor[rgb]{0.49,0.56,0.16}{{#1}}}
    \newcommand{\InformationTok}[1]{\textcolor[rgb]{0.38,0.63,0.69}{\textbf{\textit{{#1}}}}}
    \newcommand{\WarningTok}[1]{\textcolor[rgb]{0.38,0.63,0.69}{\textbf{\textit{{#1}}}}}
    \makeatletter
    \newsavebox\pandoc@box
    \newcommand*\pandocbounded[1]{%
      \sbox\pandoc@box{#1}%
      % scaling factors for width and height
      \Gscale@div\@tempa\textheight{\dimexpr\ht\pandoc@box+\dp\pandoc@box\relax}%
      \Gscale@div\@tempb\linewidth{\wd\pandoc@box}%
      % select the smaller of both
      \ifdim\@tempb\p@<\@tempa\p@
        \let\@tempa\@tempb
      \fi
      % scaling accordingly (\@tempa < 1)
      \ifdim\@tempa\p@<\p@
        \scalebox{\@tempa}{\usebox\pandoc@box}%
      % scaling not needed, use as it is
      \else
        \usebox{\pandoc@box}%
      \fi
    }
    \makeatother

    % Define a nice break command that doesn't care if a line doesn't already
    % exist.
    \def\br{\hspace*{\fill} \\* }
    % Math Jax compatibility definitions
    \def\gt{>}
    \def\lt{<}
    \let\Oldtex\TeX
    \let\Oldlatex\LaTeX
    \renewcommand{\TeX}{\textrm{\Oldtex}}
    \renewcommand{\LaTeX}{\textrm{\Oldlatex}}
    % Document parameters
    % Document title
    \title{Homework\_11}
    
    
    
    
    
    
    
% Pygments definitions
\makeatletter
\def\PY@reset{\let\PY@it=\relax \let\PY@bf=\relax%
    \let\PY@ul=\relax \let\PY@tc=\relax%
    \let\PY@bc=\relax \let\PY@ff=\relax}
\def\PY@tok#1{\csname PY@tok@#1\endcsname}
\def\PY@toks#1+{\ifx\relax#1\empty\else%
    \PY@tok{#1}\expandafter\PY@toks\fi}
\def\PY@do#1{\PY@bc{\PY@tc{\PY@ul{%
    \PY@it{\PY@bf{\PY@ff{#1}}}}}}}
\def\PY#1#2{\PY@reset\PY@toks#1+\relax+\PY@do{#2}}

\@namedef{PY@tok@w}{\def\PY@tc##1{\textcolor[rgb]{0.73,0.73,0.73}{##1}}}
\@namedef{PY@tok@c}{\let\PY@it=\textit\def\PY@tc##1{\textcolor[rgb]{0.24,0.48,0.48}{##1}}}
\@namedef{PY@tok@cp}{\def\PY@tc##1{\textcolor[rgb]{0.61,0.40,0.00}{##1}}}
\@namedef{PY@tok@k}{\let\PY@bf=\textbf\def\PY@tc##1{\textcolor[rgb]{0.00,0.50,0.00}{##1}}}
\@namedef{PY@tok@kp}{\def\PY@tc##1{\textcolor[rgb]{0.00,0.50,0.00}{##1}}}
\@namedef{PY@tok@kt}{\def\PY@tc##1{\textcolor[rgb]{0.69,0.00,0.25}{##1}}}
\@namedef{PY@tok@o}{\def\PY@tc##1{\textcolor[rgb]{0.40,0.40,0.40}{##1}}}
\@namedef{PY@tok@ow}{\let\PY@bf=\textbf\def\PY@tc##1{\textcolor[rgb]{0.67,0.13,1.00}{##1}}}
\@namedef{PY@tok@nb}{\def\PY@tc##1{\textcolor[rgb]{0.00,0.50,0.00}{##1}}}
\@namedef{PY@tok@nf}{\def\PY@tc##1{\textcolor[rgb]{0.00,0.00,1.00}{##1}}}
\@namedef{PY@tok@nc}{\let\PY@bf=\textbf\def\PY@tc##1{\textcolor[rgb]{0.00,0.00,1.00}{##1}}}
\@namedef{PY@tok@nn}{\let\PY@bf=\textbf\def\PY@tc##1{\textcolor[rgb]{0.00,0.00,1.00}{##1}}}
\@namedef{PY@tok@ne}{\let\PY@bf=\textbf\def\PY@tc##1{\textcolor[rgb]{0.80,0.25,0.22}{##1}}}
\@namedef{PY@tok@nv}{\def\PY@tc##1{\textcolor[rgb]{0.10,0.09,0.49}{##1}}}
\@namedef{PY@tok@no}{\def\PY@tc##1{\textcolor[rgb]{0.53,0.00,0.00}{##1}}}
\@namedef{PY@tok@nl}{\def\PY@tc##1{\textcolor[rgb]{0.46,0.46,0.00}{##1}}}
\@namedef{PY@tok@ni}{\let\PY@bf=\textbf\def\PY@tc##1{\textcolor[rgb]{0.44,0.44,0.44}{##1}}}
\@namedef{PY@tok@na}{\def\PY@tc##1{\textcolor[rgb]{0.41,0.47,0.13}{##1}}}
\@namedef{PY@tok@nt}{\let\PY@bf=\textbf\def\PY@tc##1{\textcolor[rgb]{0.00,0.50,0.00}{##1}}}
\@namedef{PY@tok@nd}{\def\PY@tc##1{\textcolor[rgb]{0.67,0.13,1.00}{##1}}}
\@namedef{PY@tok@s}{\def\PY@tc##1{\textcolor[rgb]{0.73,0.13,0.13}{##1}}}
\@namedef{PY@tok@sd}{\let\PY@it=\textit\def\PY@tc##1{\textcolor[rgb]{0.73,0.13,0.13}{##1}}}
\@namedef{PY@tok@si}{\let\PY@bf=\textbf\def\PY@tc##1{\textcolor[rgb]{0.64,0.35,0.47}{##1}}}
\@namedef{PY@tok@se}{\let\PY@bf=\textbf\def\PY@tc##1{\textcolor[rgb]{0.67,0.36,0.12}{##1}}}
\@namedef{PY@tok@sr}{\def\PY@tc##1{\textcolor[rgb]{0.64,0.35,0.47}{##1}}}
\@namedef{PY@tok@ss}{\def\PY@tc##1{\textcolor[rgb]{0.10,0.09,0.49}{##1}}}
\@namedef{PY@tok@sx}{\def\PY@tc##1{\textcolor[rgb]{0.00,0.50,0.00}{##1}}}
\@namedef{PY@tok@m}{\def\PY@tc##1{\textcolor[rgb]{0.40,0.40,0.40}{##1}}}
\@namedef{PY@tok@gh}{\let\PY@bf=\textbf\def\PY@tc##1{\textcolor[rgb]{0.00,0.00,0.50}{##1}}}
\@namedef{PY@tok@gu}{\let\PY@bf=\textbf\def\PY@tc##1{\textcolor[rgb]{0.50,0.00,0.50}{##1}}}
\@namedef{PY@tok@gd}{\def\PY@tc##1{\textcolor[rgb]{0.63,0.00,0.00}{##1}}}
\@namedef{PY@tok@gi}{\def\PY@tc##1{\textcolor[rgb]{0.00,0.52,0.00}{##1}}}
\@namedef{PY@tok@gr}{\def\PY@tc##1{\textcolor[rgb]{0.89,0.00,0.00}{##1}}}
\@namedef{PY@tok@ge}{\let\PY@it=\textit}
\@namedef{PY@tok@gs}{\let\PY@bf=\textbf}
\@namedef{PY@tok@ges}{\let\PY@bf=\textbf\let\PY@it=\textit}
\@namedef{PY@tok@gp}{\let\PY@bf=\textbf\def\PY@tc##1{\textcolor[rgb]{0.00,0.00,0.50}{##1}}}
\@namedef{PY@tok@go}{\def\PY@tc##1{\textcolor[rgb]{0.44,0.44,0.44}{##1}}}
\@namedef{PY@tok@gt}{\def\PY@tc##1{\textcolor[rgb]{0.00,0.27,0.87}{##1}}}
\@namedef{PY@tok@err}{\def\PY@bc##1{{\setlength{\fboxsep}{\string -\fboxrule}\fcolorbox[rgb]{1.00,0.00,0.00}{1,1,1}{\strut ##1}}}}
\@namedef{PY@tok@kc}{\let\PY@bf=\textbf\def\PY@tc##1{\textcolor[rgb]{0.00,0.50,0.00}{##1}}}
\@namedef{PY@tok@kd}{\let\PY@bf=\textbf\def\PY@tc##1{\textcolor[rgb]{0.00,0.50,0.00}{##1}}}
\@namedef{PY@tok@kn}{\let\PY@bf=\textbf\def\PY@tc##1{\textcolor[rgb]{0.00,0.50,0.00}{##1}}}
\@namedef{PY@tok@kr}{\let\PY@bf=\textbf\def\PY@tc##1{\textcolor[rgb]{0.00,0.50,0.00}{##1}}}
\@namedef{PY@tok@bp}{\def\PY@tc##1{\textcolor[rgb]{0.00,0.50,0.00}{##1}}}
\@namedef{PY@tok@fm}{\def\PY@tc##1{\textcolor[rgb]{0.00,0.00,1.00}{##1}}}
\@namedef{PY@tok@vc}{\def\PY@tc##1{\textcolor[rgb]{0.10,0.09,0.49}{##1}}}
\@namedef{PY@tok@vg}{\def\PY@tc##1{\textcolor[rgb]{0.10,0.09,0.49}{##1}}}
\@namedef{PY@tok@vi}{\def\PY@tc##1{\textcolor[rgb]{0.10,0.09,0.49}{##1}}}
\@namedef{PY@tok@vm}{\def\PY@tc##1{\textcolor[rgb]{0.10,0.09,0.49}{##1}}}
\@namedef{PY@tok@sa}{\def\PY@tc##1{\textcolor[rgb]{0.73,0.13,0.13}{##1}}}
\@namedef{PY@tok@sb}{\def\PY@tc##1{\textcolor[rgb]{0.73,0.13,0.13}{##1}}}
\@namedef{PY@tok@sc}{\def\PY@tc##1{\textcolor[rgb]{0.73,0.13,0.13}{##1}}}
\@namedef{PY@tok@dl}{\def\PY@tc##1{\textcolor[rgb]{0.73,0.13,0.13}{##1}}}
\@namedef{PY@tok@s2}{\def\PY@tc##1{\textcolor[rgb]{0.73,0.13,0.13}{##1}}}
\@namedef{PY@tok@sh}{\def\PY@tc##1{\textcolor[rgb]{0.73,0.13,0.13}{##1}}}
\@namedef{PY@tok@s1}{\def\PY@tc##1{\textcolor[rgb]{0.73,0.13,0.13}{##1}}}
\@namedef{PY@tok@mb}{\def\PY@tc##1{\textcolor[rgb]{0.40,0.40,0.40}{##1}}}
\@namedef{PY@tok@mf}{\def\PY@tc##1{\textcolor[rgb]{0.40,0.40,0.40}{##1}}}
\@namedef{PY@tok@mh}{\def\PY@tc##1{\textcolor[rgb]{0.40,0.40,0.40}{##1}}}
\@namedef{PY@tok@mi}{\def\PY@tc##1{\textcolor[rgb]{0.40,0.40,0.40}{##1}}}
\@namedef{PY@tok@il}{\def\PY@tc##1{\textcolor[rgb]{0.40,0.40,0.40}{##1}}}
\@namedef{PY@tok@mo}{\def\PY@tc##1{\textcolor[rgb]{0.40,0.40,0.40}{##1}}}
\@namedef{PY@tok@ch}{\let\PY@it=\textit\def\PY@tc##1{\textcolor[rgb]{0.24,0.48,0.48}{##1}}}
\@namedef{PY@tok@cm}{\let\PY@it=\textit\def\PY@tc##1{\textcolor[rgb]{0.24,0.48,0.48}{##1}}}
\@namedef{PY@tok@cpf}{\let\PY@it=\textit\def\PY@tc##1{\textcolor[rgb]{0.24,0.48,0.48}{##1}}}
\@namedef{PY@tok@c1}{\let\PY@it=\textit\def\PY@tc##1{\textcolor[rgb]{0.24,0.48,0.48}{##1}}}
\@namedef{PY@tok@cs}{\let\PY@it=\textit\def\PY@tc##1{\textcolor[rgb]{0.24,0.48,0.48}{##1}}}

\def\PYZbs{\char`\\}
\def\PYZus{\char`\_}
\def\PYZob{\char`\{}
\def\PYZcb{\char`\}}
\def\PYZca{\char`\^}
\def\PYZam{\char`\&}
\def\PYZlt{\char`\<}
\def\PYZgt{\char`\>}
\def\PYZsh{\char`\#}
\def\PYZpc{\char`\%}
\def\PYZdl{\char`\$}
\def\PYZhy{\char`\-}
\def\PYZsq{\char`\'}
\def\PYZdq{\char`\"}
\def\PYZti{\char`\~}
% for compatibility with earlier versions
\def\PYZat{@}
\def\PYZlb{[}
\def\PYZrb{]}
\makeatother


    % For linebreaks inside Verbatim environment from package fancyvrb.
    \makeatletter
        \newbox\Wrappedcontinuationbox
        \newbox\Wrappedvisiblespacebox
        \newcommand*\Wrappedvisiblespace {\textcolor{red}{\textvisiblespace}}
        \newcommand*\Wrappedcontinuationsymbol {\textcolor{red}{\llap{\tiny$\m@th\hookrightarrow$}}}
        \newcommand*\Wrappedcontinuationindent {3ex }
        \newcommand*\Wrappedafterbreak {\kern\Wrappedcontinuationindent\copy\Wrappedcontinuationbox}
        % Take advantage of the already applied Pygments mark-up to insert
        % potential linebreaks for TeX processing.
        %        {, <, #, %, $, ' and ": go to next line.
        %        _, }, ^, &, >, - and ~: stay at end of broken line.
        % Use of \textquotesingle for straight quote.
        \newcommand*\Wrappedbreaksatspecials {%
            \def\PYGZus{\discretionary{\char`\_}{\Wrappedafterbreak}{\char`\_}}%
            \def\PYGZob{\discretionary{}{\Wrappedafterbreak\char`\{}{\char`\{}}%
            \def\PYGZcb{\discretionary{\char`\}}{\Wrappedafterbreak}{\char`\}}}%
            \def\PYGZca{\discretionary{\char`\^}{\Wrappedafterbreak}{\char`\^}}%
            \def\PYGZam{\discretionary{\char`\&}{\Wrappedafterbreak}{\char`\&}}%
            \def\PYGZlt{\discretionary{}{\Wrappedafterbreak\char`\<}{\char`\<}}%
            \def\PYGZgt{\discretionary{\char`\>}{\Wrappedafterbreak}{\char`\>}}%
            \def\PYGZsh{\discretionary{}{\Wrappedafterbreak\char`\#}{\char`\#}}%
            \def\PYGZpc{\discretionary{}{\Wrappedafterbreak\char`\%}{\char`\%}}%
            \def\PYGZdl{\discretionary{}{\Wrappedafterbreak\char`\$}{\char`\$}}%
            \def\PYGZhy{\discretionary{\char`\-}{\Wrappedafterbreak}{\char`\-}}%
            \def\PYGZsq{\discretionary{}{\Wrappedafterbreak\textquotesingle}{\textquotesingle}}%
            \def\PYGZdq{\discretionary{}{\Wrappedafterbreak\char`\"}{\char`\"}}%
            \def\PYGZti{\discretionary{\char`\~}{\Wrappedafterbreak}{\char`\~}}%
        }
        % Some characters . , ; ? ! / are not pygmentized.
        % This macro makes them "active" and they will insert potential linebreaks
        \newcommand*\Wrappedbreaksatpunct {%
            \lccode`\~`\.\lowercase{\def~}{\discretionary{\hbox{\char`\.}}{\Wrappedafterbreak}{\hbox{\char`\.}}}%
            \lccode`\~`\,\lowercase{\def~}{\discretionary{\hbox{\char`\,}}{\Wrappedafterbreak}{\hbox{\char`\,}}}%
            \lccode`\~`\;\lowercase{\def~}{\discretionary{\hbox{\char`\;}}{\Wrappedafterbreak}{\hbox{\char`\;}}}%
            \lccode`\~`\:\lowercase{\def~}{\discretionary{\hbox{\char`\:}}{\Wrappedafterbreak}{\hbox{\char`\:}}}%
            \lccode`\~`\?\lowercase{\def~}{\discretionary{\hbox{\char`\?}}{\Wrappedafterbreak}{\hbox{\char`\?}}}%
            \lccode`\~`\!\lowercase{\def~}{\discretionary{\hbox{\char`\!}}{\Wrappedafterbreak}{\hbox{\char`\!}}}%
            \lccode`\~`\/\lowercase{\def~}{\discretionary{\hbox{\char`\/}}{\Wrappedafterbreak}{\hbox{\char`\/}}}%
            \catcode`\.\active
            \catcode`\,\active
            \catcode`\;\active
            \catcode`\:\active
            \catcode`\?\active
            \catcode`\!\active
            \catcode`\/\active
            \lccode`\~`\~
        }
    \makeatother

    \let\OriginalVerbatim=\Verbatim
    \makeatletter
    \renewcommand{\Verbatim}[1][1]{%
        %\parskip\z@skip
        \sbox\Wrappedcontinuationbox {\Wrappedcontinuationsymbol}%
        \sbox\Wrappedvisiblespacebox {\FV@SetupFont\Wrappedvisiblespace}%
        \def\FancyVerbFormatLine ##1{\hsize\linewidth
            \vtop{\raggedright\hyphenpenalty\z@\exhyphenpenalty\z@
                \doublehyphendemerits\z@\finalhyphendemerits\z@
                \strut ##1\strut}%
        }%
        % If the linebreak is at a space, the latter will be displayed as visible
        % space at end of first line, and a continuation symbol starts next line.
        % Stretch/shrink are however usually zero for typewriter font.
        \def\FV@Space {%
            \nobreak\hskip\z@ plus\fontdimen3\font minus\fontdimen4\font
            \discretionary{\copy\Wrappedvisiblespacebox}{\Wrappedafterbreak}
            {\kern\fontdimen2\font}%
        }%

        % Allow breaks at special characters using \PYG... macros.
        \Wrappedbreaksatspecials
        % Breaks at punctuation characters . , ; ? ! and / need catcode=\active
        \OriginalVerbatim[#1,codes*=\Wrappedbreaksatpunct]%
    }
    \makeatother

    % Exact colors from NB
    \definecolor{incolor}{HTML}{303F9F}
    \definecolor{outcolor}{HTML}{D84315}
    \definecolor{cellborder}{HTML}{CFCFCF}
    \definecolor{cellbackground}{HTML}{F7F7F7}

    % prompt
    \makeatletter
    \newcommand{\boxspacing}{\kern\kvtcb@left@rule\kern\kvtcb@boxsep}
    \makeatother
    \newcommand{\prompt}[4]{
        {\ttfamily\llap{{\color{#2}[#3]:\hspace{3pt}#4}}\vspace{-\baselineskip}}
    }
    

    
    % Prevent overflowing lines due to hard-to-break entities
    \sloppy
    % Setup hyperref package
    \hypersetup{
      breaklinks=true,  % so long urls are correctly broken across lines
      colorlinks=true,
      urlcolor=urlcolor,
      linkcolor=linkcolor,
      citecolor=citecolor,
      }
    % Slightly bigger margins than the latex defaults
    
    \geometry{verbose,tmargin=1in,bmargin=1in,lmargin=1in,rmargin=1in}
    
    

\begin{document}
    
    \maketitle
    
    

    
    Probability for Data Science

UC Berkeley, Spring 2026

Ani Adhikari

CC BY-NC-SA 4.0

This content is protected and may not be shared, uploaded, or
distributed.

    \hypertarget{homework-11-due-monday-april-20th-at-5-pm}{%
\section{Homework 11 (Due Monday, April 20th at 5
PM)}\label{homework-11-due-monday-april-20th-at-5-pm}}

    \hypertarget{instructions}{%
\subsubsection{Instructions}\label{instructions}}

\hypertarget{questions-1-2-and-3-can-be-answered-based-on-tuesdays-lecture-and-past-content-alone.}{%
\paragraph{\texorpdfstring{\textbf{Questions 1, 2, and 3 can be answered
based on Tuesday's lecture (and past content)
alone.}}{Questions 1, 2, and 3 can be answered based on Tuesday's lecture (and past content) alone.}}\label{questions-1-2-and-3-can-be-answered-based-on-tuesdays-lecture-and-past-content-alone.}}

\hypertarget{we-strongly-encourage-you-to-attempt-questions-1-5-and-6-independently-relying-on-the-provided-textbook-sections-and-your-own-reasoning.-if-you-find-yourself-stuck-we-recommend-attending-office-hours-or-homework-party-or-posting-your-questions-on-ed.-if-youre-having-trouble-getting-started-on-problems-without-the-use-of-ai-we-encourage-you-to-follow-our-recommendation.}{%
\paragraph{\texorpdfstring{\textbf{We strongly encourage you to attempt
Questions 1, 5, and 6 independently, relying on the provided textbook
sections and your own reasoning. If you find yourself stuck, we
recommend attending office hours or homework party or posting your
questions on Ed. If you're having trouble getting started on problems
without the use of AI, we encourage you to follow our
recommendation.}}{We strongly encourage you to attempt Questions 1, 5, and 6 independently, relying on the provided textbook sections and your own reasoning. If you find yourself stuck, we recommend attending office hours or homework party or posting your questions on Ed. If you're having trouble getting started on problems without the use of AI, we encourage you to follow our recommendation.}}\label{we-strongly-encourage-you-to-attempt-questions-1-5-and-6-independently-relying-on-the-provided-textbook-sections-and-your-own-reasoning.-if-you-find-yourself-stuck-we-recommend-attending-office-hours-or-homework-party-or-posting-your-questions-on-ed.-if-youre-having-trouble-getting-started-on-problems-without-the-use-of-ai-we-encourage-you-to-follow-our-recommendation.}}

Your homeworks will generally have two components: a written portion and
a portion that also involves code. Written work should be completed on
paper, and coding questions should be done in the notebook. Start the
work for the written portions of each section on a new page. You are
welcome to \(\LaTeX\) your answers to the written portions, but staff
will not be able to assist you with \(\LaTeX\) related issues.

It is your responsibility to ensure that both components of the homework
are submitted completely and properly to Gradescope. \textbf{Make sure
to assign each page of your pdf to the correct question. Refer to the
bottom of the notebook for submission instructions.}

    \hypertarget{how-to-do-your-homework}{%
\subsubsection{How to Do Your Homework}\label{how-to-do-your-homework}}

The point of homework is for you to try your hand at using what you've
learned in class. The steps to follow:

\begin{itemize}
\tightlist
\item
  Go to lecture and sections, and also go over the relevant text
  sections before starting on the homework. This will remind you what
  was covered in class, and the text will typically contain examples not
  covered in lecture. The weekly Study Guide will list what you should
  read.
\item
  Work on some of the practice problems before starting on the homework.
\item
  Attempt the homework problems by yourself with the text, section work,
  and practice materials all at hand. Sometimes the week's lab will help
  as well. The two steps above will help this step go faster and be more
  fruitful.
\item
  At this point, seek help if you need it. Don't ask how to do the
  problem --- ask how to get started, or for a nudge to get you past
  where you are stuck. Always say what you have already tried. That
  helps us help you more effectively.
\item
  For a good measure of your understanding, keep track of the fraction
  of the homework you can do by yourself or with minimal help. It's a
  better measure than your homework score, and only you can measure it.
\end{itemize}

    \hypertarget{rules-for-homework}{%
\subsubsection{Rules for Homework}\label{rules-for-homework}}

\begin{itemize}
\tightlist
\item
  Every answer should contain a calculation or reasoning. For example, a
  calculation such as \((1/3)(0.8) + (2/3)(0.7)\) or
  \texttt{sum({[}(1/3)*0.8,\ (2/3)*0.7{]})}is fine without further
  explanation or simplification. If we want you to simplify, we'll ask
  you to. But just \({5 \choose 2}\) by itself is not fine; write ``we
  want any 2 out of the 5 frogs and they can appear in any order'' or
  whatever reasoning you used. Reasoning can be brief and abbreviated,
  e.g.~``product rule'' or ``not mutually exclusive.''
\item
  You may consult others (see ``How to Do Your Homework'' above) but you
  must write up your own answers using your own words, notation, and
  sequence of steps.
\item
  We'll be using Gradescope. You must submit the homework according to
  the instructions at the end of homework set.
\end{itemize}

\hypertarget{we-will-not-grade-assignments-which-do-not-have-pages-correctly-selected-for-each-question.}{%
\subsection{\texorpdfstring{\textbf{We will not grade assignments which
do not have pages correctly selected for each
question.}}{We will not grade assignments which do not have pages correctly selected for each question.}}\label{we-will-not-grade-assignments-which-do-not-have-pages-correctly-selected-for-each-question.}}

    \newpage

    \hypertarget{waiting-for-a-random-coin-to-land-heads}{%
\subsection{1. Waiting for a Random Coin to Land
Heads}\label{waiting-for-a-random-coin-to-land-heads}}

Let \(X\) be a random proportion. Given \(X=p\), let \(T\) be the number
of tosses till a \(p\)-coin lands heads.

\textbf{a)} Let \(P(X = 1/10) = 1/4\), \(P(X = 1/7) = 1/4\), and
\(P(X = 1/3) = 1/2\). Find \(E(T)\). Start with \(E(T|X = p)\). Read
through
\href{https://data140.org/textbook/content/chapter-09/expectation-by-conditioning/}{Ch.
9.2} (Expecatation by Conditioning) and
\href{https://data140.org/textbook/content/chapter-08/expectations-of-functions/}{Ch.
8.3} (Expectations of Functions) for a refresher on expectations of
discrete distributions and expectation by conditioning.

\begin{align*}
	E[T] &= E[E[T \mid X]]  \\ 
	&= E[T \mid X = \frac{1}{10}]P(X = \frac{1}{10}) + E[T \mid X = \frac{1}{7}]P(X =\frac{1}{7}) + E[T \mid X = \frac{1}{3}]P(X=\frac{1}{3}) \\
&= 10 \cdot \frac{1}{4} + 7 \cdot \frac{1}{4} + 3 \cdot \frac{1}{3} \\
&= \frac{17}{4} + 1 \\
&= 3 + \frac{1}{4} \\
\end{align*}



\textbf{b)} Find \(E(T)\) if \(X\) has the beta \((r, s)\) density for
some \(r > 1\). Simplify all integrals and Gamma functions in your
answer. Remember that \(\Gamma(w+1) = w\Gamma(w)\). Use the same setup
as part (a) and skim
\href{https://data140.org/textbook/content/chapter-15/expectation/}{Ch.
15.3} (Expectation of Continuous Random Variables) for help on finding
the expectation.

\begin{align*}
	E[T] &= E[E[T \mid X]] \\
	&= \int_{0}^{1} E[T \mid X = x] P(X = x)dx \\
	&= \int_{0}^{1} \frac{1}{x} B(r,s) x^{r-1}(1-x)^{s-1}dx \\
	&= B(r,s)\int_{0}^{1} x^{(r-1)-1}(1-x)^{s-1}dx \\
	&= B(r,s) \cdot \frac{1}{B(r-1, s)} \\
	&= \frac{\Gamma(r+s)}{\Gamma(r)\Gamma(s)} \cdot \frac{\Gamma(r-1)\Gamma(s)}{\Gamma(r-1+s)} \\
	&= \frac{(r+s-1)\Gamma(r+s-1)}{(r-1)\Gamma(r-1)\Gamma(s)} \cdot \frac{\Gamma(r-1)\Gamma(s)}{\Gamma(r-1+s)} \\
	&= \frac{r+s-1}{r-1} \\
\end{align*}




\textbf{c)} Let \(X\) have the beta \((r, s)\) density. For fixed
\(k > 0\), find the posterior density of \(X\) given \(T = k\). Identify
it as one of the famous ones and state its name and parameters. This is
a straightforward application of Ch.
\href{https://data140.org/textbook/content/chapter-21/updating-and-prediction/\#updating-the-posterior-distribution-of-x-given-s-n}{21.1.3}
(Beta Updates).



\begin{align*}
	f_{X \mid T = k}(p) &= \frac{f_{X}(p)dp P(T=k)}{P(T=k)} \\
	&\propto C(r,s)p^{r-1}(1-p)^{s-1} \cdot p^{1}(1-p)^{k-1} \\
			    &\propto p^{(r+1)-1}(1-p)^{(s+k-1) -1} \\
\end{align*}


This is the Beta(r+1, s+k-1) density. 






    \newpage

    \hypertarget{discrete-and-continuous-random-selections}{%
\subsection{2. Discrete and Continuous Random
Selections}\label{discrete-and-continuous-random-selections}}

Fix a positive integer \(n\), and let \(p\) be strictly between 0 and 1.

Suppose Person A picks a number uniformly in the interval \((0, n)\).
Let \(X\) be the number Person A picks.

Suppose that independently of Person A, Person B picks an integer \(Y\)
according to the binomial \((n, p)\) distribution, for example by using
\texttt{stats.binom.rvs(n,\ p,\ size=1)} or by tossing a \(p\)-coin
\(n\) times and recording the number of heads.

Find \(P(X < Y)\).


$X \sim $ uniform(0,n), $Y \sim \Binom(n, p)$ number of heads. 

$X = \{0,1,2,3,4,5,\dots n\}$. $Y = \{0,1,2,3,4,5,\dots, n\}$ 


This is the bottom right triangle of the n by n cube. 



\begin{align*}
	P(X < Y) &= \int_{0}^{n} P(Y=y)\int_{0}^{y} P(X=x) dxdy \\
	&= \sum_{y=0}^{n} P(Y=y) \int_{0}^{y} P(X=x) dx \\
&=  	\sum_{y=0}^{n} \binom{n}{y} p^{y}(1-p)^{n-y} \int_{0}^{y}P(X=x)dx\\
&= \sum_{y=0}^{n} \binom{n}{y}p^{y}(1-p)^{n-y} \int_{0}^{y}\frac{1}{n}dx \\
&= \sum_{y=0}^{n} \binom{n}{y} p^{y}(1-p)^{n-y} \left[ \frac{1}{n}x \right]_{0}^{y} \\
&= \sum_{y=0}^{n} \binom{n}{y} p^{y}(1-p)^{n-y} \left( \frac{y}{n} \right)  \\
&= \sum_{y=0}^{n} \frac{n!}{y!(n-y)!} \cdot \frac{y}{n} p^{y}(1-p)^{n-y} \\
&= \sum_{y=0}^{n}\frac{(n-1)!}{(y-1)!(n-y)!} p^{y}(1-p)^{n-y} \\
&= \sum_{y=0}^{n} \frac{(n-1)!}{(y-1)!( (n-1) - (y-1))!} p^{y}(1-p)^{n-y}  \\
&= \sum_{y=0}^{n} p \binom{n-1}{y-1} p^{y-1} (1-p)^{n-y} \\
&= p \cdot 1 \\
&= p \\
\end{align*}
    \newpage

    \hypertarget{prediction-with-sums}{%
\subsection{3. Prediction with Sums}\label{prediction-with-sums}}

Let \(X_1, X_2, \ldots, X_n\) be i.i.d. with expectation \(\mu\) and
variance \(\sigma^2\). Let \(S = \sum_{i=1}^n X_i\).

\textbf{a)} Find the least squares predictor of \(S\) based on \(X_1\),
and find the mean squared error (MSE) of the predictor.

\begin{align*}
	E[S \mid X_{1}] &= E[X_1 + X_2 + X_3 + \dots X_{n} \mid X_{1}] \\
	&= E[X_1 \mid X_1] + E[X_2 \mid X_1] + E[X_3 \mid X_1] + \dots + E[X_{n} \mid X_1 ]  \\
	&= X_1 + (n-1)E[X_2 \mid X_1], \quad \text{by symmetry}\\ 
	&= X_1 + (n-1)E[X_2], \quad \text{by independence} \\
	&= X_{1} + (n-1)\mu  \\
\end{align*}

\begin{align*}
	MSE(E[S \mid X_{1}]) &= E\left[(S  - E[S \mid X_{1}])^{2}\right] \\
	&= E((X_1 + X_2 + \dots + X_n - (X_1 + (n-1)\mu ))^{2}) \\
	&= E[(X_2 + X_3 + \dots + X_{n} - (n-1)\mu )^{2}] \\ 
	&= \Var(\sum_{i=2}^{n} X_{i}) \\
	&= \Var(X_2) + \Var(X_3) + \dots + \Var(X_{n}), \text{by independence} \\
	&= (n-1)\Var(X_{2}), \quad \text{by symmetry} \\
	&= (n-1)\sigma^{2} \\
\end{align*}


\textbf{b)} Find the least squares predictor of \(X_1\) based on \(S\),
and find the MSE of the predictor. Is the predictor a linear function of
\(S\)? If so, it must also be the best among all linear predictors based
on \(S\), which is commonly known as the regression predictor.

{[}Consider whether your predictor in (b) would be different if \(X_1\)
were replaced by \(X_2\), or by \(X_3\), or by \(X_i\) for any fixed
\(i\). Then use symmetry and the additivity of conditional
expectation.{]}


\[
E[X_1 \mid S] = E[S - (X_2 + X_3 + \dots + X_{n})  \mid S] = S - (n-1)E[X_1 \mid S]
\]
\begin{align*}
	E[X_1 \mid S] &= S - (n-1)E[X_{1} \mid S] \\
	nE[X_1 \mid S] &= S\\
	&= \frac{S}{n} \\
	&= \frac{1}{n} \sum_{i=1}^{n} X_{n} \\
	&= \bar X \\
\end{align*}

Yes, the function is a linear function of $S$.
    \newpage

    \hypertarget{sections}{%
\subsection{4. Sections}\label{sections}}

A class of 60 students has three sections. Summary statistics for scores
on Quiz 7: - Section 1: 25 students, mean 25, SD 3 - Section 2: 20
students, mean 23, SD 2 - Section 3: 15 students, mean 27, SD 4

Let \(S\) be the Quiz 7 score of a student picked at random from the
class. Use the code cell below to calculate numerical answers for parts
a and b.

\textbf{a)} Find \(E(S)\).
Let $I$ indicate the section where $\Omega = \{1,2,3\}$, with probability $\{\frac{25}{60}, \frac{20}{60}, \frac{15}{60}\}$ 


\begin{align*}
	E[S] &= E[E[S \mid I]] \\ 
	&= \frac{25}{60} \cdot 25 + \frac{20}{60} \cdot 23 + \frac{15}{60} \cdot 27 \\
	&= \frac{149}{6} \\
	&= 24.333 \\
\end{align*}


\textbf{b)} Find \(SD(S)\).

\[
	\Var(S) = E[\Var(S \mid I)] + \Var(E[S \mid I]) \\
\]

\begin{align*}
	E[\Var(S \mid I)] &= \frac{25}{60} \cdot 9 + \frac{20}{60} \cdot 4 + \frac{15}{60} \cdot 16 \\
	&= \frac{109}{12} \\
	&= 9.0833 \\
\end{align*} 

\begin{align*}
	\Var(E[S \mid I]) &= E[S^{2} \mid I] - E[S^{2} \mid I]^{2} \\
	&= \frac{25}{60} \cdot 25^{2} + \frac{20}{60} 23^{2} + \frac{15}{60} 27^{2} - \left( \frac{149}{6} \right) ^{2} \\
	&= \frac{242}{9} \\
	&= 26.89 \\
\end{align*}

\[
\Var(S) = \frac{109}{12} + \frac{242}{9} = \frac{1295}{36} = 35.97
\]

\[
SD(S) = \sqrt{\Var(S))} \approx 6
\]
 C  \newpage

    \hypertarget{two-colored-die}{%
\subsection{5. Two-Colored Die}\label{two-colored-die}}

This problem is an exercise in conditioning. To find expectation, read
through
\href{https://data140.org/textbook/content/chapter-09/expectation-by-conditioning/}{Ch.
9.2} (Expectation by Conditioning). To find variance, read through
\href{https://data140.org/textbook/content/chapter-22/variance-by-conditioning/}{Ch.
22.3} (Variance by Conditioning) and note the decomposition of variance.
For examples of computing expectation and variance using conditioning,
skim through
\href{https://data140.org/textbook/content/chapter-22/examples/}{Ch.
22.4} (Examples). Problems similar to this one are Ch. 22 Ex.1, which is
done in Wednesdays's section, and Ch. 22 Ex.3, which is done in Friday's
section.

A die that has two green faces and four blue faces is rolled repeatedly.
Let \(R\) be the number of rolls till both colors have appeared.

In each part below, find the numerical value.

\textbf{a)} Find \(E(R)\).

Let $I_0$ be the first color that appears. $I_{0} = 1$ if green, and $0$ otherwise.  

\begin{align*}
	E[R] &= E[R \mid I_0] \\
	&= 1 + \frac{1}{\frac{2}{6}} \cdot \frac{4}{6} + \frac{1}{\frac{4}{6}}\cdot  \frac{2}{6} \\
	&= 1 +  \frac{6}{2} \cdot \frac{4}{6} + \frac{6}{4} \cdot \frac{2}{6}\\
	&= 1 + 2 + \frac{1}{2} \\
	&= 3 + \frac{1}{2} \\
\end{align*}






\textbf{b)} Find \(SD(R)\).



\begin{align*}
	\Var(R) &= E[\Var(R \mid I_0)] + \Var(E[R \mid I_{0}])\\ 
\end{align*}


\begin{align*}
	E[\Var(R \mid I_0)] &= \Var(R \mid I_0 = 1)P(I_0 = 1) + \Var(R \mid I_0 = 0)P(I_1=0) \\
	&= \frac{2}{6} / \frac{4}{6}^{2} \cdot \frac{2}{6} + \frac{4}{6} / \frac{2}{6}^{2} \cdot \frac{4}{6} \\
	&= \frac{4}{16} + \frac{16}{4} \\
	&= \frac{1}{4} + 4 \\
	&= 4.25 \\
\end{align*}



\begin{align*}
	\Var(E[R \mid I_0]) &= E[R^{2} \mid I_0] - E[R \mid I_0]^{2} \\
	&= (1 + \frac{6}{2})^{2} \frac{4}{6} + (1 + \frac{6}{4})^{2}\frac{2}{6}- \frac{7}{2}^{2}\\
	&= 4^{2} \cdot \frac{4}{6} + 2.5^{2} \cdot \frac{2}{6} - \frac{7}{2}^{2} \\
	&= 0.5 \\
\end{align*}

\[
\Var(R) = 4.25  + 0.5 = 4.75
\]

\[
	SD(R) = \sqrt{4.75} = 2.18
\]
    \newpage

    \hypertarget{accident-claims}{%
\subsection{6. Accident Claims}\label{accident-claims}}

Before you do this exercise, carefully study the results of
\href{https://data140.org/textbook/content/chapter-22/examples/\#random-sum}{Ch.
22.4.3} (Random Sum). You will have to follow similar steps. For
additional information on the results of this chapter, skim through
\href{https://data140.org/textbook/content/chapter-22/variance-by-conditioning/}{Ch.
22.3} (Variance by Conditioning) and note the decomposition of variance.
A problem similar to this one is Ch. 22 Ex.4, which is done in Friday's
section.

Poisson processes are standard models used by actuaries and industrial
engineers.

At a street intersection, accidents happen according to a Poisson
process at a rate of \(\lambda\) per month. Each accident leads to an
insurance claim of a random dollar amount. Assume that the amounts of
the claims are i.i.d. with mean \(\mu\) and SD \(\sigma\), independent
of the process of accidents.

Find the expectation and variance of the total amount claimed for
accidents at that intersection in a \textbf{year}.


The total number of accidents in a year jis a Poisson Process $\{N(t), t \ge 0\}$ with rate $12\lambda$.

Let $T$ be the total amount claimed for accidents at that intersection. 


\begin{align*}
	E[T] &= E[E[T \mid N]] \\
	&= \sum_{n=0}^{\infty }E[T \mid N=n]P(N=n) \\
	&= \sum_{0}^{\infty } n\mu e^{-12\lambda} \frac{(12\lambda)^{n}}{n!} \\
	&= \mu  \cdot E[N] \\
	&= \mu \cdot 12\lambda \\
	&= 12\lambda\mu  \\
\end{align*}

\begin{align*}
	\Var(T) &= E[\Var(T \mid N)] + \Var(E[T \mid N]) \\
\end{align*}

\begin{align*}
	E[\Var(T \mid N)] &= \sum_{n=0}^{\infty } n\sigma^{2} P(N=n) \\
	&= \sigma^{2} \sum_{n=0}^{\infty }nP(N=n) \\
	&= \sigma^{2} E[N] \\
	&= \sigma^{2} \cdot 12\lambda \\
	&= 12\lambda\sigma^{2} \\
\end{align*}

\begin{align*}
	\Var(E[T \mid N]) &= \sum_{n=0}^{\infty } \mu ^{2} n P(N=n) \\
	&= 12\lambda\mu ^{2} \\
\end{align*}

\[
\Var(T) =  12\lambda(\sigma^{2} + \mu ^{2}) 
\]

\[
	SD(T) = \sqrt{12\lambda(\sigma^{2} + \mu ^{2})}
\]



    \newpage

    \hypertarget{submission-instructions}{%
\subsection{Submission Instructions}\label{submission-instructions}}

Many assignments throughout the course will have a written portion and a
code portion. Please follow the directions below to properly submit both
portions.

\hypertarget{written-portion}{%
\subsubsection{Written Portion}\label{written-portion}}

\begin{itemize}
\tightlist
\item
  Scan all the pages into a PDF. You can use any scanner or a phone
  using applications such as CamScanner. Please \textbf{DO NOT} simply
  take pictures using your phone.
\item
  Please start a new page for each question. If you have already written
  multiple questions on the same page, you can crop the image in
  CamScanner or fold your page over (the old-fashioned way). This helps
  expedite grading.
\item
  It is your responsibility to check that all the work on all the
  scanned pages is legible.
\item
  If you used \(\LaTeX\) to do the written portions, you do not need to
  do any scanning; you can just download the whole notebook as a PDF via
  LaTeX.
\end{itemize}

\hypertarget{code-portion}{%
\subsubsection{Code Portion}\label{code-portion}}

\begin{itemize}
\tightlist
\item
  Save your notebook using
  \texttt{File\ \textgreater{}\ Save\ and\ Checkpoint}.
\item
  Generate a PDF file using
  \texttt{File\ \textgreater{}\ Download\ As\ \textgreater{}\ PDF\ via\ LaTeX}.
  This might take a few seconds and will automatically download a PDF
  version of this notebook.

  \begin{itemize}
  \tightlist
  \item
    If you have issues, please post a follow-up on the general Homework
    11 Ed thread.
  \end{itemize}
\end{itemize}

\hypertarget{submitting}{%
\subsubsection{Submitting}\label{submitting}}

\begin{itemize}
\tightlist
\item
  Combine the PDFs from the written and code portions into one PDF.
  \href{https://smallpdf.com/merge-pdf}{Here} is a useful tool for doing
  so.
\item
  Submit the assignment to Homework 11 on Gradescope.
\item
  \textbf{Make sure to assign each page of your pdf to the correct
  question.}
\item
  \textbf{It is your responsibility to verify that all of your work
  shows up in your final PDF submission.}
\end{itemize}

If you are having difficulties scanning, uploading, or submitting your
work, please read the
\href{https://edstem.org/us/courses/94054/discussion/7533569}{Ed Thread}
on this topic and post a follow-up on the general Homework 11 Ed thread.


    % Add a bibliography block to the postdoc
    
    
    
\end{document}
